\RequirePackage{xr-hyper}
\documentclass[pdflatex,sn-mathphys-num]{sn-jnl}
\geometry{a4paper,left=22mm,right=22mm,top=22mm,bottom=24mm}
\usepackage{graphicx,amsmath,amssymb,booktabs,array,xcolor,microtype,longtable}
\usepackage[version=4]{mhchem}
\usepackage{placeins,hyperref}
\definecolor{revisionblue}{RGB}{0,75,160}
\newif\ifreviewcopy
\ifdefined\CleanCopy\reviewcopyfalse\else\reviewcopytrue\fi
\newcommand{\rev}[2]{\ifreviewcopy{\color{revisionblue}#2\textsuperscript{\scriptsize[R1-#1]}}\else#2\fi}
\renewcommand{\thefigure}{S\arabic{figure}}
\renewcommand{\thetable}{S\arabic{table}}
\newcommand{\file}[1]{\texttt{\detokenize{#1}}}
\providecommand{\sitab}[1]{Supplementary Table~\ref{#1}}
\providecommand{\sifig}[1]{Supplementary Fig.~\ref{#1}}
\hypersetup{hidelinks}
\externaldocument[M-]{revision1/sn-article-final}
% Use standard floats so externally scaled tabulars are not reboxed by threeparttable.
\makeatletter
\let\table\tableorg
\let\endtable\endtableorg
\makeatother
\raggedbottom
\begin{document}
\title[Supplementary Information]{Supplementary Information: When Does a Foundation Model Help? Process-Specific Reliability of Machine-Learned Interatomic Potentials}
\author[1]{\fnm{Yi} \sur{Cao}}
\author[1]{\fnm{Peter} \sur{Mastracco}}
\author*[1]{\fnm{Paulette} \sur{Clancy}}\email{pclancy3@jhu.edu}
\affil[1]{\orgdiv{Department of Chemical and Biomolecular Engineering},\orgname{Johns Hopkins University},\orgaddress{\city{Baltimore},\state{Maryland},\postcode{21218},\country{USA}}}
\maketitle

\section{Reference data, training routes and pointwise evaluation}\label{sec:reference}
\subsection{Reference scope and computational provenance}
The reference archive contains Quantum ESPRESSO calculations using PBE-based electronic structure for Cr-doped \ce{Sb2Te3}~\citep{giannozzi2009quantum,perdew1996generalized}. The parent AIMD methods describe D3 dispersion, SG15 pseudopotentials, 100/400~Ry wavefunction/charge-density cutoffs, cold smearing of 0.0147~Ry, an electronic threshold of $10^{-6}$~Ry, and 300, 600 and 1200~K sampling. These are parent-protocol metadata, not verified common settings for every subsequently combined migration and sliding record. The configuration-specific input and status take precedence over that general description. In particular, a path extracted from an archived profile is not promoted to a converged DFT minimum-energy path by attaching a peak value to it.

Three reference collections must be distinguished. The legacy five-image paths G1, G2 and Q1--Q3 support the original grouped overlap and retraining controls. The canonical 81-atom, ten-image paths supply the selected main migration figures, with in-gap, gap-to-QL and within-QL reference peaks of 1.5243, 0 and 1.8449~eV. Their source status retains SCF warnings. Finally, the fixed sliding reference profiles have peaks of 1.6538~eV for the pristine system and 2.5708~eV for Cr in the gap. We retain their profiles and endpoint-return checks separately. Neither the legacy G1 reference of 0.336~eV nor its errors can be substituted for the canonical in-gap reference.

\subsection{Model matrix and training-data comparisons}\label{sec:training}
The contemporary pointwise evaluation contains Scratch-600K, FT-600K, Scratch-MultiT and FT-MultiT, each with training seeds 123, 234 and 345. The 600~K pair uses the same 1,786 training configurations and 311/267 validation/test configurations. The MultiT pair has 2,659 training and 767 test configurations in the pointwise release. Thus the plot's pooled 5,358/801 and 7,977/2,301 energy records are three predictions per unique train/test configuration, not larger independent datasets.

Fine-tuned models belong to the grouped-A2, OMAT-FT-start, stage-two compiled lineage. Scratch models in the contemporary evaluation use the registered $L_{\max}=1$ architecture. The original four-model and later pathway campaigns contain other checkpoints and are not assumed to be architecture-matched. Model-file and data-file SHA256 identities are retained in the accompanying registry. This study compares the practical routes represented by those implementations; it does not isolate pretraining from architecture, loss design and optimization in a controlled factorial experiment.

\begin{figure}[htbp]
\centering
\includegraphics[width=\textwidth]{revision1/figures/main/SI/method-P02_palette_updated_pre_vector_export.pdf}
\caption{\textbf{Training routes and energy prediction.} \textbf{a,} Schematic mapping from atomic coordinates through descriptors and a neural energy model to an extensive total energy, expressed as a sum of atomic contributions. \textbf{b--d,} Foundation pretraining on a broad corpus, specialist training from random initialization on target data, and adaptation of a foundation checkpoint on target data. The drawing explains the relationship between routes; it does not specify an identical architecture or training budget for the archived checkpoints, nor does it depict all equivariant message-passing operations in MACE. Implementation-specific model identities accompany the numerical data.}
\label{fig:training_design}
\end{figure}

\begin{table}[htbp]
\centering\small
\caption{\textbf{Current uncentred held-out errors.} Energy RMSE in meV atom$^{-1}$ and Cartesian force-component RMSE in meV \AA$^{-1}$, reported as means and sample standard deviations over training seeds 123, 234 and 345. The underlying 24 train/test evaluations are retained in the current-release source table. Temperature conditions have different test compositions, so this is not a common-test comparison between 600~K and MultiT.}
\label{tab:current_pointwise}
\input{revision1/tables/current_pointwise.tex}
\end{table}

The parity visualization subtracts a DFT trajectory-cluster mean from both reference and prediction energies. Shared train/test clusters use one pooled reference mean; a cluster occurring only in the test split uses its test reference mean. This is a display transformation, not a calibration available to the deployed model. Residuals are unchanged by subtracting the same offset from both axes, but correlations after cluster-centering describe fluctuations within those clusters. The uncentred errors in Table~\ref{tab:current_pointwise} are therefore reported independently of the plotted correlation. The test Scratch-MultiT force display excludes 106 components outside $\pm10$~eV \AA$^{-1}$, while the numerical error calculation retains them.

\subsection{Force-label audit}\label{sec:force_audit}
The exporter issue identified in review concerns a real distinction between missing and zero forces. Missing QE force labels must never be converted to an all-zero array and admitted to force-supervised training. The corrected policy is fail-closed: a force-bearing training export requires validated force data; energy-only geometry exports must not impersonate force-labelled frames.

The retained audit covers 9,163 records across historical FT-600K, historical FT-MultiT, grouped FT-600K and grouped Scratch-5\% train/validation/test files. Every record has a force field. Nine records have exactly zero forces, all explicitly labelled one-atom Cr, Sb or Te \file{IsolatedAtom} energy-reference configurations. There are no unclassified or NEB-derived zero-force records in those audited files. The historical training counts are 1,891 and 3,015 for FT-600K and FT-MultiT, respectively; they are distinct from the current grouped datasets. The manifest records frame identity, source headers, force-array presence, zero-force classification, source-file checksums and split membership. The audit does not establish the provenance of a different checkpoint merely because it has the same family label.

\subsection{Grouped splitting, overlap and data-volume controls}\label{sec:overlap}
The grouped protocol keeps complete source trajectories within training, validation or test. It prevents adjacent frames from the same trajectory being distributed across partitions, but does not imply that physically similar configurations from different trajectories cannot occur. The original frame-shuffle baseline is not used as the independent generalization estimate.

The overlap audit compares 25 legacy migration images with the full 600~K grouped training file and the 89-configuration Scratch-5\% subset. It finds zero exact global matches, zero exact Cr-local matches, and no SOAP cosine distance at or below $10^{-4}$. Minimum SOAP distances are $5.4\times10^{-4}$ and $6.4\times10^{-4}$, respectively. The minimum comparable Cr-local RMSD is 2.06~\AA, above the 0.05~\AA\ near-duplicate threshold. Species-compatible RMSD comparisons cover only two training candidates for the full file and one for the small subset; SOAP covers all audited images. These results rule out the defined duplicates within the audited files, rather than proving a large distributional separation.

The legacy G1 grouped fixed-image errors are $0.47\pm0.12$~eV for FT-600K, $1.20\pm0.27$~eV for full Scratch and $0.68\pm0.20$~eV for Scratch-5\%. The original 0.16~eV point estimate is not reproduced. The low-data subset contains 89 training configurations, not 1,000, and does not match the 1,786-frame FT training budget. We use it only to describe data-volume sensitivity. The full-data scratch and fine-tuned runs provide the target-data-matched comparison.

\begin{table}[htbp]
\centering\small
\caption{\textbf{Legacy grouped control, kept distinct from the current canonical paths.} Values are mean $\pm$ sample standard deviation across three training seeds, except for the single MPA0 checkpoint. G1 is a 61-atom, five-image reference with a QE forward peak of 0.336~eV and an unconverged path status. The last column gives diagnostic fixed-image discrepancies, not validated migration-barrier errors. The current pointwise release in Table~\ref{tab:current_pointwise} supplies freshly recomputed static metrics; minor differences from the rounded legacy table do not alter the separate G1 control.}
\label{tab:a2_grouped_audit_summary}
\begin{tabular}{lrrr}
\toprule
Route & Train frames & G1 discrepancy (eV) & Training seeds\\
\midrule
FT-600K & 1,786 & $0.47\pm0.12$ & 3\\
Scratch & 1,786 & $1.20\pm0.27$ & 3\\
Scratch-5\% & 89 & $0.68\pm0.20$ & 3\\
MPA0 & zero-shot & 0.24 & checkpoint\\
\bottomrule
\end{tabular}
\end{table}

\begin{figure}[htbp]
\centering
\includegraphics[width=\textwidth]{revision1/figures/grouped_retraining_overlap.pdf}
\caption{\textbf{Legacy grouped retraining and train--path overlap control.} \textbf{a,} Absolute fixed-image discrepancies on G1 for FT-600K, Scratch and Scratch-5\%, retaining individual training runs and their mean $\pm$ sample standard deviation. The foundation symbol is a single checkpoint, not a training-seed average. The dashed 0.16~eV line is the superseded frame-shuffle point estimate. \textbf{b,} Minimum training-set SOAP and composition-compatible local RMSD distances normalized by their respective near-duplicate thresholds. Values greater than one lie outside the declared thresholds. G1 and the 25-image audit belong to the legacy reference collection; the 0.336~eV G1 peak is not the 1.5243~eV canonical in-gap reference used in the main seed matrix.}
\label{fig:grouped_overlap}
\end{figure}

\begin{figure}[htbp]
\centering
\includegraphics[width=.92\textwidth]{revision1/figures/task_conditioned/fig_task_conditioned_panel_a.pdf}
\caption{\textbf{Configuration-space context for the legacy grouped reference partition.} SOAP-derived principal components are fitted on the grouped training partition and used to visualize the 1,786 training, 311 validation and 267 test configurations together with the audited migration images. Symbols identify provenance and split. The projection complements the explicit nearest-neighbour check; visible proximity or separation alone is not a duplicate criterion or a physical error estimate.}
\label{fig:grouped_configuration_space}
\end{figure}

\begin{figure}[htbp]
\centering
\includegraphics[width=\textwidth]{revision1/figures/task_conditioned/fig_original_training_rmse_four_model.pdf}
\caption{\textbf{Training histories for the original four-model baseline.} Validation error histories are reconstructed from 87 archived default-head checkpoints. Each series retains its recorded training chronology. The curves document the original optimization histories; they are not a common-test, matched-epoch or independent-trajectory generalization comparison, and they do not replace the current grouped pointwise evaluation.}
\label{fig:legacy_training_history}
\end{figure}
\FloatBarrier

\section{Migration references, operators and seed-resolved discrepancies}\label{sec:migration}
\subsection{Metric definitions}
For image set $\Gamma=\{X_0,\ldots,X_n\}$, let
\begin{align}
\Delta E_M(i)&=E_M(X_i)-E_M(X_0),\\
B_M^\rightarrow(\Gamma)&=\max_i\Delta E_M(i),\\
B_M^\leftarrow(\Gamma)&=\max_i E_M(X_i)-E_M(X_n).
\end{align}
Including the initial image makes the forward peak nonnegative. A monotonically downhill profile consequently has $B^\rightarrow=0$, irrespective of its endpoint energy or reverse peak. A peak discrepancy can therefore be small while the profile is inaccurate. Whole-band spans, forward peaks and reverse peaks are distinct quantities.

Fixed scoring evaluates identical geometries for every calculator. Endpoint-initialized, self-relaxed NEB instead evaluates an optimized image set $\Gamma_M$. Its optimizer convergence describes the target-model calculation; it does not certify that $\Gamma_M$ is the DFT path or even connects the same basins after relaxation. The main figures use the finite diagnostic values as supplied and do not convert them into rate constants.

\subsection{Current main-figure matrix and status reconciliation}\label{sec:matrix_status}
The main seed matrix contains 15 model rows and eight tasks, giving 120 finite numerical records. Absolute discrepancies were independently recomputed as the absolute difference between the predicted peak and the task reference; all 120 match the supplied values. The three migration references are 1.524299, 0 and 1.8449~eV. The two sliding references are 1.653834726 and 2.5708301675~eV. These numerical definitions are used consistently in the figure-associated table.

The originally supplied matrix labelled 75 records as \file{FINITE_DIAGNOSTIC}, 42 as \file{CONVERGED} and three as \file{NONCONVERGED}. A final join against the canonical Scratch-600K raw campaign matches all nine self-relaxed records by seed, path and model SHA256. It identifies two incorrect convergence labels: seed 123/within-QL and seed 234/in-gap are nonconverged, with residuals 0.2748 and 0.3613~eV \AA$^{-1}$. The first record also has an incorrect zero peak in the supplied matrix; its raw forward peak is 2.3507668938545976~eV. The second peak, 0.3116233847104013~eV, is unchanged. Main Fig.~\ref{M-fig:taskmatrix} corrects that value and both labels and marks all five nonconverged records with crosses. The self-relaxed panel limit is extended to retain the corrected point. No finite value is omitted. The revised tags are 75 finite-diagnostic, 40 optimizer-converged and five nonconverged entries; optimizer convergence is not DFT validation. The as-supplied table, original PDF, hash-matched nine-record comparison and explicit correction audit are retained alongside the corrected data.

There is also a lineage distinction. Fig.~\ref{M-fig:profiles} retains an earlier four-condition profile ensemble: its scratch curve is legacy, its foundation band pools multiple checkpoints, and some fine-tuned fixed/self curves use different available members. Fig.~\ref{M-fig:taskmatrix} is a later seed-expanded matrix. Family labels must not be used to claim identical checkpoint pairing across those releases.

\begin{table}[htbp]
\centering\scriptsize
\caption{\textbf{Archived membership behind the selected profile figure.} Counts are model/checkpoint members for each path class in main Fig.~\ref{M-fig:profiles}, as recorded in the R02 coverage table. The repeated counts do not turn different foundation checkpoints into training-seed replicates. FT-MultiT has no accepted within-QL self-relaxed curve. Contemporary Scratch-600K additions from R05 belong to a later release and are not assigned to the selected four-condition figure.}
\label{tab:profile_coverage}
\begin{tabular}{lrrrr}
\toprule
Operator / path & Foundation & Legacy Scratch & FT-600K & FT-MultiT\\
\midrule
Fixed / in gap & 9 & 3 & 3 & 2\\
Fixed / gap to QL & 9 & 3 & 3 & 2\\
Fixed / within QL & 9 & 3 & 3 & 2\\
Self / in gap & 9 & 3 & 3 & 3\\
Self / gap to QL & 9 & 3 & 3 & 3\\
Self / within QL & 9 & 3 & 3 & 0\\
\bottomrule
\end{tabular}
\end{table}

\begin{table}[htbp]
\centering\small
\caption{\textbf{Current fixed-path reference discrepancies.} Absolute differences from the supplied task peak, in eV. Fine-tuned and scratch rows report mean $\pm$ sample standard deviation over the three training seeds; foundation rows are individual checkpoints. These statistics are computed from the 120-record Fig.~\ref{M-fig:taskmatrix} table. Migration references retain SCF warnings, and sliding values do not replace a full-profile or endpoint-return check.}
\label{tab:current_neb_errors}
\resizebox{\textwidth}{!}{\input{revision1/tables/current_neb_errors.tex}}
\end{table}

\subsection{Legacy image sets and corrected earlier estimates}\label{sec:legacy_paths}
The original 61-atom, one-Cr G1 path contains five images. Its QE forward image-set peak is 0.336~eV, or 0.34~eV at two decimals. The earlier 0.41~eV foundation result came from an 82-atom, two-Cr, twelve-image record, with inconsistent endpoint/internal calculators and a band-serialization defect. It is withdrawn. Uniform rescoring of that preserved band gives a forward peak of 0, reverse peak of 0.38 and whole-band span of 1.04~eV, none interchangeable with the 0.41~eV claim. We do not interpret the difference as relaxation improving the same path.

\begin{table}[htbp]
\centering\scriptsize
\caption{\textbf{Legacy QE path references and convergence states.} Forward and reverse peak values are relative to the corresponding endpoints. Unconverged paths supply image-set energy diagnostics only. A converged path label refers to its archived path calculation and does not certify an unrelated model-relaxed band. These paths are distinct from the canonical main-figure set.}
\label{tab:dft_neb_barriers_revised}
\resizebox{\textwidth}{!}{\input{revision1/tables/dft_neb_barriers.tex}}
\end{table}

\begin{table}[htbp]
\centering\scriptsize
\caption{\textbf{Legacy migration path registry.} Class and geometry-based naming connect the retained reference figures and path-resolved controls to their released identifiers. The registry prevents a five-image G/Q reference being confused with a later ten-image canonical path of the same broad class.}
\label{tab:path_registry}
\resizebox{\textwidth}{!}{\input{revision1/tables/path_registry.tex}}
\end{table}
\input{revision1/tables/scratch5_path_resolved_descriptive.tex}

The Q1 Scratch discrepancy of $0.51\pm0.11$~eV is stable across the three legacy training seeds, whereas Q2 and Q3 have larger variation. The corresponding FT-600K discrepancies remain approximately 10~eV. These observations replace the earlier attribution to chance. They describe reference-image discrepancies on the stated paths, not universal extrapolation performance or converged kinetic errors.

\begin{figure}[htbp]
\centering
\includegraphics[width=.48\textwidth]{revision1/figure_pages/assets/fig_041.pdf}
\includegraphics[width=.48\textwidth]{revision1/figure_pages/assets/fig_042.pdf}\\
\includegraphics[width=.48\textwidth]{revision1/figure_pages/assets/fig_043.pdf}
\includegraphics[width=.48\textwidth]{revision1/figure_pages/assets/fig_044.pdf}\\
\includegraphics[width=.48\textwidth]{revision1/figure_pages/assets/fig_045.pdf}
\caption{\textbf{Retained fixed-image profiles for the legacy path collection.} \textbf{a--e,} Q1, Q2, G1, Q3 and G2, respectively, evaluated on each path's supplied QE geometries. Model energies are endpoint-referenced within each panel. Q paths approach or enter a quintuple layer; G paths remain in the gap under the source classification. The archived QE paths are marked unconverged, so agreement tests common-image energy assignments rather than a validated activation energy. These profiles support the legacy grouped controls and are not relabelled as the canonical main-figure paths.}
\label{fig:legacy_profiles}
\end{figure}

\begin{figure}[htbp]
\centering
\includegraphics[width=\textwidth]{revision1/figures/figS_dft_neb_mep_profiles.pdf}
\caption{\textbf{Additional archived within-QL QE profiles.} The panels retain the L1 and C1--C5 image-set profiles. Their endpoint-referenced forward and reverse peaks and archived convergence labels are listed in Table~\ref{tab:dft_neb_barriers_revised}. Path identity and the electronic/path convergence record, rather than profile appearance, determine whether an activation-energy interpretation is supported.}
\label{fig:legacy_qe_profiles}
\end{figure}

\begin{figure}[htbp]
\centering
\includegraphics[width=.49\textwidth]{revision1/figures/foundation_grouped/path_1-4_energy_profiles_corrected.pdf}
\includegraphics[width=.49\textwidth]{revision1/figures/foundation_grouped/path_1-4_energy_profiles_family_aggregated.pdf}
\caption{\textbf{Foundation-checkpoint sensitivity on the legacy in-gap image set.} \textbf{a,} Ten checkpoint profiles on the same five G1 geometries, with relative energies referenced to the first image. \textbf{b,} Family-aggregated profiles, with variation across available submodels. These bands describe checkpoint variation, not independent training-seed uncertainty. The QE G1 forward reference is 0.336~eV and the path is unconverged. This control therefore tests fixed-image calibration; it does not supply a physical migration-barrier ranking or the canonical 1.5243~eV reference.}
\label{fig:legacy_foundation}
\end{figure}

\begin{figure}[htbp]
\centering
\includegraphics[width=\textwidth]{revision1/figures/figS_neb_operators.pdf}
\caption{\textbf{Why the two original NEB evaluations were not comparable.} \textbf{a,} Fixed-image energy assignments for the 61-atom, one-Cr G1 set, with its 0.336~eV QE reference. \textbf{b,} Archived model-driven optimization behaviour, illustrating that the band can become unstable when forces move the images. The separate 82-atom, two-Cr record used for the withdrawn 0.41~eV claim is not a matched version of G1. Neither an unstable band nor a differently defined band span is used as a physical barrier.}
\label{fig:legacy_operator_audit}
\end{figure}

\subsection{Additional dopant diagnostics}
The supplied Cu, Li and Na profiles extend the qualitative fixed-image comparison to other dopants in the same host. They do not supply the complete grouped training, MD and path-validation controls available for the primary Cr study. The screened QE path table also retains electronic-convergence limitations. We therefore use the figure as an additional diagnostic, rather than as validation of transfer across materials.

\begin{figure}[htbp]
\centering
\includegraphics[width=.92\textwidth]{revision1/figures/main/SI/fig_NEB-fixed-path_Cu-Li-Na-Cr.pdf}
\caption{\textbf{Fixed-path foundation-model diagnostics for four dopants in \ce{Sb2Te3}.} Rows show Cr (\textbf{a--c}), Cu (\textbf{d--f}), Li (\textbf{g--i}) and Na (\textbf{j--l}); columns show in-gap, gap-to-QL and within-QL configurations. Curves compare the displayed MACE-MP, MACE-OMAT-0 and MACE-MH checkpoint groups with the screened QE path table. Energies are relative to the initial image, and the horizontal coordinate is normalized image index. Group bands describe available checkpoint variation. A downhill reference can have a zero forward peak while retaining a nonzero reverse peak or endpoint difference. The QE records retain convergence limitations, so the curves are diagnostic energy assignments rather than validated activation energies. They do not establish the cross-material generality of the main model-selection conclusions.}
\label{fig:other_dopants}
\end{figure}
\FloatBarrier

\section{Collective sliding}\label{sec:sliding}
Sliding profiles compare a prescribed relative displacement of the layers with endpoint-initialized, model-relaxed bands for pristine and Cr-containing structures. In the fixed view, all models are intended to score the supplied geometries; in the self-relaxed view, the intermediate geometries depend on the model. The first-image energy is subtracted separately for each profile.

The nominal endpoint-return check is important because a profile can approach a reference peak while ending at a different relative energy. Before attributing that discrepancy to interaction range or forgetting, the identical periodic configuration, species/atom mapping, cell, calculator head and energy reference must be verified. The supplied plots do not isolate the cause of an endpoint mismatch.

\begin{figure}[htbp]
\centering
\includegraphics[width=\textwidth]{revision1/figures/main/SI/fig_NEB_interlayer-sliding.pdf}
\caption{\textbf{Sliding profiles across foundation, fine-tuned and scratch routes.} \textbf{a,b,} Fixed-path relative energies for pristine and Cr-in-gap \ce{Sb2Te3}, respectively. \textbf{c,d,} Endpoint-initialized self-relaxed profiles for the same system categories. Solid coloured summaries distinguish the displayed training conditions and their available member variation; the dashed dark curve is the supplied DFT reference. The normalized coordinate labels image order, rather than a guaranteed common arc length after relaxation. DFT reference peaks are 1.6538~eV for pristine and 2.5708~eV for Cr-in-gap sliding. Peak agreement, profile agreement, endpoint return and self-relaxation convergence are separate requirements. Self-relaxed curves are MLIP pathway diagnostics and have not been DFT-rescored as minimum-energy paths.}
\label{fig:sliding_all}
\end{figure}

\begin{figure}[htbp]
\centering
\includegraphics[width=\textwidth]{revision1/figures/main/SI/fig_NEB_interlayer-sliding-foundation-model.pdf}
\caption{\textbf{Foundation-family sensitivity of interlayer sliding.} \textbf{a,b,} Fixed-path relative energies for pristine and Cr-in-gap structures. \textbf{c,d,} Corresponding self-relaxed pathway diagnostics. MACE-MP, MACE-OMAT-0 and MACE-MH are displayed as checkpoint groups, with pointwise member variation; shaded bands are not training-seed or DFT uncertainty. Each curve is referenced to its own first image. The DFT profile is retained as a supplied reference, and no rate or validated activation barrier is inferred from the self-relaxed comparison. The failure to recover a nominally equivalent endpoint motivates an input/calculator consistency check before a physical-mechanism attribution.}
\label{fig:sliding_foundation}
\end{figure}
\FloatBarrier

\section{Matched MD protocol and trajectory horizons}\label{sec:md}
The common continuation starts from an archived state at step 200,000 with finite preserved momenta and a 2,050-atom periodic cell. The target is step 1,000,000. The timestep is 1~fs, the thermostat target is 600~K, the Langevin friction parameter in the runner is 0.01, and the trajectory output interval is 200 steps. The continuation manifest records MD seed 1100; it is not a separate training seed. Thus a run ending at 236~ps has 36~ps of continuation from the common state, not a complete 1~ns observation.

All curves retain their recorded end times. Threshold-reached runs are not extended by interpolation, replication or zero padding. The three foundation targets have provisional target-reached source status; reaching a step count is not independent validation of their physical dynamics. For the purpose of the reported result, the directly supported comparison is their observed longer trajectory horizon.

\begin{table}[htbp]
\centering\scriptsize
\caption{\textbf{Observed MD horizons for every displayed model.} Clock time is reconstructed from the last step using the 1~fs timestep. Continuation time subtracts the common 200~ps starting clock. Threshold-reached runs retain their finite observations but are not counted as completed continuations. Foundation entries describe distinct checkpoints, while trained-family entries describe three training seeds; they are not independent starting-state/velocity replicas.}
\label{tab:md_horizons}
\resizebox{\textwidth}{!}{\input{revision1/tables/current_md_horizons.tex}}
\end{table}

Kinetic temperature is computed from archived momenta. Hydrostatic pressure is $-\mathrm{tr}(\sigma)/3$ using the saved stress tensor. In-plane MSD uses fractional-coordinate unwrapping in the triclinic cell, and velocity autocorrelation uses multiple time origins with normalization by its zero-lag value. These observables use the supported time window of each run.

Pair coordination numbers integrate the partial RDF to fixed first-minimum cutoffs of 3.2~\AA\ (Cr--Sb), 4.1~\AA\ (Cr--Te) and 4.2~\AA\ (Sb--Te). The all-atom RDF uses ordered pair counts, $r_{\max}=10$~\AA, bins of 0.1~\AA\ and every twentieth eligible frame. Its denominator accumulates $N(N-1)/V$ times shell volume. Condition summaries receive display-only smoothing; raw pair-count outputs remain unchanged. Unequal horizons mean the structural summaries do not represent an identical stationary observation window across all models.

The original asymmetric short FT-600K comparison is not used. No Green--Kubo heat-current record supports a thermal-conductivity estimate for the present campaign. We do not recycle earlier diffusion coefficients into this figure, nor derive common long-time transport coefficients from failure-bounded traces.
\section{Current trajectory-feature analysis}\label{sec:representation}
\subsection{Feature definition and limitations}
The selected Fig.~\ref{M-fig:representation} contains 927 sampled frames across 15 models, with 4--101 frames per model. Each legacy feature vector has one energy value, 6,150 Cartesian force components, a 30-bin neighbour-distance histogram, ten SOAP PCA entries and one all-zero sensitivity placeholder. The current audit removes the final column before computing any distance or new projection, leaving 6,191 columns. A constant zero coordinate adds nothing to Euclidean or cosine distances; it is not a calculated force response.

Three methodological limits govern interpretation. Model energies use different reference offsets. Features have different units and scales. In addition, the ten SOAP PCA summaries were fitted separately within each trajectory, so their axes do not define one common physical coordinate basis. These output features must not be described as internal MACE activations. A common 30-bin neighbour histogram provides a separate, structurally defined basis for a limited original-space comparison.

\begin{table}[htbp]
\centering\small
\caption{\textbf{Original-space metrics for the current 927-frame dataset.} Silhouette labels are the 15 model identities, using the same sampled rows as the main figure. The output-feature vector excludes the inactive placeholder. Standardization is column-wise z-scoring over the combined diagnostic sample and is not a training transformation. Neighbour histograms use common distance bins, avoiding the separate SOAP PCA bases. Scores are descriptive and are affected by unequal trajectory populations.}
\label{tab:current_original_space}
\resizebox{\textwidth}{!}{\input{revision1/tables/current_original_space.tex}}
\end{table}

\subsection{Projection sensitivity}
The displayed $t$-SNE uses perplexity 30 and random state 42. PHATE uses random state 42 and its historical default neighbour setting; the source run reported disconnected components. UMAP uses 15 neighbours, minimum distance 0.1, Euclidean distance and random/transform seeds 42. Panel c colours a 45-bin two-dimensional histogram count with a maximum count of 26; it is not a free energy or sampling probability estimate.

The current audit repeats $t$-SNE directly on the 6,191-column matrix at perplexities 5, 30 and 50, with nominal seeds 0, 42 and 123 and PCA initialization. Nominal seeds are algorithm settings, not independent physics replicates. The resulting silhouette range is 0.079--0.133. This sensitivity and the scale dependence of original-space scores preclude a claim that the visual clusters identify disjoint learned mechanisms.

\begin{table}[htbp]
\centering\small
\caption{\textbf{Current $t$-SNE sensitivity checks.} Scores use the 15 model labels. Trustworthiness uses five neighbours and compares each projection with the same unprojected input matrix. These metrics assess a visualization and do not measure model prediction uncertainty or DFT accuracy.}
\label{tab:current_tsne_sensitivity}
\input{revision1/tables/current_tsne_sensitivity.tex}
\end{table}

\begin{figure}[htbp]
\centering
\includegraphics[width=\textwidth]{revision1/figures/figS_current_feature_audit.pdf}
\caption{\textbf{Scale and projection checks on the same current dataset.} \textbf{a,} Silhouette scores before projection for raw and standardized output features and for the common neighbour histogram. The inactive legacy column is removed. \textbf{b,} $t$-SNE silhouette values across three perplexities and three nominal algorithm seeds. Original-space values and per-run settings appear in Tables~\ref{tab:current_original_space} and \ref{tab:current_tsne_sensitivity}. The raw output-space score is 0.063, the z-scored value is $-0.059$, and the common histogram score is 0.056. These results support a descriptive population diagnostic, not a claim of native latent-space separation or a validated acquisition criterion.}
\label{fig:current_feature_audit}
\end{figure}
\FloatBarrier

\section{Common-configuration error surrogates and SHAP}\label{sec:shap}
The current analysis replaces the earlier 150-frame, three-model SHAP reconstruction. It uses the shared 267-frame, 86-atom test trajectory \file{2L_o3_t3_600K_aimd_1}, identified by explicit frame alignment and a common reference-energy check. All four model families and all three training seeds are evaluated on these same configurations. The additional 500 MultiT test frames are excluded from this cross-family analysis because corresponding 600~K predictions are unavailable; they remain in the source tables.

The target is the absolute total-energy error divided by atom count, in meV atom$^{-1}$. Inputs comprise 390 periodic, inner-averaged DScribe SOAP features for Cr/Sb/Te, with $r_{\rm cut}=5$~\AA, $n_{\max}=4$ and $l_{\max}=4$, plus 17 explicitly computed structural features. SOAP power-spectrum labels identify species channels and radial/angular indices. They are not MACE neurons, independently manipulable bonds or evidence that an MLIP has internalized a specified physical interaction.

For the final TreeSHAP fit, each model--seed case uses a gradient-boosting regressor with 200 estimators, learning rate 0.05, maximum depth 8, subsampling fraction 0.8 and random state 42. A separate lightweight diagnostic uses 50 estimators, depth 4 and a shuffled five-fold KFold split with random state 42. Its cross-validation metrics are given below. Frames come from one trajectory, so shuffled-fold accuracy does not establish performance on an independent trajectory. The final full-sample fit has approximately unit in-sample $R^2$ and is not itself an out-of-sample validation.

\begin{table}[htbp]
\centering\small
\caption{\textbf{Current error-surrogate fidelity diagnostic.} Means and sample standard deviations over the three underlying MLIP training seeds. Each value summarizes the separate five-fold diagnostic estimator on the common 267-frame sample. The MAE unit is meV atom$^{-1}$, rather than the total-energy eV unit used in the superseded 150-frame analysis. These values do not quantify an independent-trajectory generalization error for the final TreeSHAP estimator.}
\label{tab:current_shap_fidelity}
\input{revision1/tables/current_shap_fidelity.tex}
\end{table}

The main figure uses one shared global top-20 feature ordering in both panels. Bars summarize mean absolute SHAP and sample standard deviation over the three training seeds within each condition; the heat map retains all 12 model--seed columns. Large between-seed variation limits any single-feature explanation. The current ranking does not retain the earlier assertion that a specific Cr--Cr channel explains migration fidelity.

The accompanying executable workflow takes the common XYZ and per-frame evaluation tables as inputs, builds the 407 computed features, checks frame/reference alignment, fits the surrogates and writes the per-seed metrics and attribution tables. The feature audit and figure-specific source ledger are also included. These files replace the empty-stub API as the reproducibility route for the results in this manuscript.
\FloatBarrier

\input{revision1/appendix/endpoint_perturbation_appendix.tex}
\FloatBarrier
\bibliography{revision1/references-camera-ready}
\end{document}
